problem:  
Let \( (M^n, g) \) be a complete noncompact Riemannian manifold satisfying the **borderline Ricci curvature decay**
\[
\mathrm{Ric}(x) \ge -\frac{C}{(1+r(x))^2}, \qquad r(x) = d(x,x_0),
\]
and assume **noncollapsed volume growth**
\[
\operatorname{Vol}(B(x_0,r)) \ge c\,r^n \quad \text{for all } r\ge1,
\]
for some constants \(c, C>0\).  

For each \(p>2\), consider the space of global weak \(p\)-harmonic functions \(u \in W^{1,p}_{\mathrm{loc}}(M)\) satisfying the nonlinear equation
\[
\operatorname{div}(|\nabla u|^{p-2}\nabla u) = 0
\]
and the **sublinear growth** condition
\[
|u(x)| = o(r(x)) \qquad \text{as } r(x)\to\infty.
\]

**Question (Nonlinear Liouville Rigidity at the Borderline Curvature Decay):**  
Under these assumptions, is it true that for each fixed \(p>2\), every sublinear \(p\)-harmonic function on \(M\) must be constant?  

Equivalently, does the curvature condition \(\mathrm{Ric}\ge -C(1+r)^{-2}\) already guarantee global \(p\)-parabolicity (in the sense that all \(p\)-energy finite functions with vanishing flux at infinity are constant) whenever the volume growth is noncollapsed and Euclidean to first order?  

A further refinement asks whether, in the absence of such nonlinear Liouville rigidity, the existence of a nontrivial sublinear \(p\)-harmonic function must correspond to **nonuniqueness of tangent cones at infinity** in the measured Gromov–Hausdorff limit.

Why is it a "good" problem:  
This version isolates a **sharp nonlinear generalization** of an unresolved threshold in geometric analysis—the behavior of harmonic functions under borderline curvature decay—and extends it to the highly nonlinear, degenerate elliptic regime of \(p\)-harmonic analysis.  

For \(\mathrm{Ric}\ge0\), classical results (Li–Yau, Holopainen, Kotschwar–Ni) show that sublinear \(p\)-harmonic functions are necessarily constant, and manifolds with noncollapsing volume growth exhibit \(p\)-parabolicity. However, when curvature merely satisfies \(\mathrm{Ric}\ge -C(1+r)^{-2}\), the standard gradient estimates and Sobolev inequalities destabilize exactly at the critical scale, and existing techniques do not provide either Liouville or non-Liouville outcomes.  

This problem therefore focuses on the **borderline between parabolicity and hyperbolicity** for nonlinear potential theory on manifolds of almost nonnegative curvature, while coupling it to the global geometry at infinity. A positive resolution would link \(p\)-energy growth, curvature decay, and asymptotic cone rigidity under minimal assumptions, suggesting a unified nonlinear analogue of the Colding–Minicozzi–Cheeger–Colding harmonic structure theory.  

Its formulation is simple—“Does borderline curvature decay still enforce a Liouville theorem for all \(p>2\)?”—yet its solution demands new techniques in nonlinear analysis, metric measure theory, and geometric PDE.  It is thus positioned squarely at a deep and uncharted intersection of nonlinear elliptic theory and asymptotic geometry, meeting all criteria of a “good” research-level problem.